\documentclass[11pt]{article}
\usepackage[a4paper,margin=25mm]{geometry}
\usepackage{amsmath,amssymb}
\usepackage[hidelinks]{hyperref}
\title{Restricted order of asymptotic bases\\
\large Lean proofs for the original Erd\H{o}s Problem \#338 questions}
\author{Leo Gao\\\small AI-assisted research and formalization}
\date{27 November 2026}
\newcommand{\N}{\mathbb N}
\newcommand{\RR}{R}
\setlength{\emergencystretch}{3em}
\begin{document}
\maketitle
\begin{abstract}
For any ordinary basis of order $h$, a rank-two reflected-intersection
test characterizes equality with the restricted order. Its local
inputs are ordinary good-element pair sums and a bounded conflict
graph of fringe labels. The Lean package checks this criterion, uniform boundedness, and
Chen--Yang--Zhao's existence and both finite-deletion results.
The proof of the sharp growth rate is on paper. The kernel checks
the stated theorems, not priority or the adequacy of the proposed
equality characterization.
\end{abstract}

\section{The four questions}
We use at most $h$ positive summands. Write $S_h(A)$ when repetition is
allowed and $R_h(A)$ when summands must be distinct. The ordinary
and restricted orders are the least positive $h$ for which the
corresponding sumset contains all sufficiently large integers.
The four questions are those in concluding remark 3 on p.~206 of
\href{https://www.renyi.hu/~p_erdos/1980-22.pdf\#page=6}{Erd\H{o}s and Graham (1980)}.

\begin{enumerate}
\item \textbf{Existence.} A finite ordinary-order basis has finite
restricted order exactly when distinct finite subset sums meet
every residue class for every modulus.
This is Theorem 2 of
\href{https://arxiv.org/html/2611.33191v1}{Chen, Yang and Zhao}.
The supplied manuscript also derives the canonical variant:
the eventual tail gcd is $g$, and it suffices to check subset sums
of the finitely many elements not divisible by $g$, modulo $g$.
\item \textbf{Uniform bound.} For ordinary order at most $h$, finite
restricted order is bounded by a function of $h$ alone.
The companion \emph{Sharp boundedness} proof gives
$F(h)\le \exp(O(2^h h\log(2h)))$, and a mixed-radix construction
gives $\log\log F(h)=h\log2+O(\log h)$.
The Lean proof supplies a separate, coarser explicit bound.
A non-strict bound suffices for the original strict formulation by adding one.
\item \textbf{Equality.} The rank-two criterion below is necessary
and sufficient for arbitrary bases. It eliminates the use of
higher-rank restricted representation functions.
\item \textbf{Finite deletions.} Both variants are affirmative by
Chen--Yang--Zhao, Theorem 1: if deleting any finite set preserves
being an ordinary basis of some finite order, finite restricted
order follows. Requiring the same ordinary order is stronger.
\end{enumerate}

\section{The proposed equality criterion}
For $h\ge2$, let $B$ be the elements of $A$ centered in at most
$h-2$ nontrivial symmetric pairs from $A$, and put $G=A\setminus B$.
For a weight $s$, list the possible nonempty sets of fringe labels
in a distinct representation by at most two terms: the singleton
$\{s\}$ if $s\in B$, the singletons $\{b\}$ with $b\in B$ and
$s-b\in G$, and the unordered pairs $\{b,s-b\}\subseteq B$
of distinct terms. Denote the family by $M(s)$.
It is a matching: no two of its members share any label.
Declare $s$ immune if $s\in S_2(G)$ or $|M(s)|\ge h-1$.
Every remaining weight is weak if $M(s)$ is nonempty, and is
unavailable otherwise.

Write $h=2m+e$, $e\in\{0,1\}$. Decompose a target into $m$
pair weights and, when $e=1$, an optional ordinary singleton.
Require no unavailable weight. Choose one fringe support per
weak occurrence, with all these supports disjoint from each other
and from the singleton if it is in $B$. There is no choice to record
at immune weights. For each decomposition this is a finite
conflict-graph test with at most $m(h-2)$ vertices.

\paragraph{Why it works.}
A distinct sum immediately gives an accepted decomposition by
grouping its terms in pairs. Conversely realize the weak weights
first. An immune weight can be realized while avoiding all used
fringe labels: there are at most $h-2$ such labels, which cannot
hit $h-1$ disjoint alternatives. If the weight is in $S_2(G)$
and its ordinary good representation repeats a midpoint, an unused
symmetric pair achieves the same sum. After gluing, only good
summands can repeat. Replace a repeated good midpoint by a symmetric
pair avoiding the surviving summands, increasing the number of
distinct labels at every step. The process ends with a distinct
sum of rank at most $h$.

Thus the tested targets are exactly $R_h(A)$. If $S_h(A)$ is
cofinite, it suffices to apply the test eventually on
$2\cdot B+S_{h-2}(A)$, because any failed target with an ordinary
representation must repeat a fringe element in that representation.
Order one is immediate.

\section{Why this is a candidate answer, and its limits}
The criterion is not a restatement using $R_h$ or its
representation count as input. It uses only rank-two reflected
intersections, ordinary sums, and compatibility of a uniformly
bounded number of fringe labels at each decomposition. It discards
all endpoint data at immune weights. It gives a necessary-and-sufficient condition in
terms of lower-rank data, rather than taking the full restricted
sumset as an input.

The criterion still has an eventual quantifier over targets.
For a positive all-fringe basis with unique two-term sums and no
singleton--pair coincidences, every occupied positive weight has one
weak witness. The test then keeps the actual endpoints and checks
disjointness. It does not classify the global distribution of
obstructions in arbitrary bases. An exact equivalence can be
kernel-checked; whether the criterion is satisfactory for the
historical question is a distinct judgment. No acceptance,
peer review, or priority is claimed here. The accompanying
proof map identifies the five checked declarations and the
separate paper proof of the sharp bound.

\section{Files for review}
The full LaTeX proofs of the equality criterion, its more general
blocker formulation, and the quantitative growth theorem are
included. The source bundle also includes a checked Lean 4 project,
a fixed dependency revision, a build script, and an axiom audit.
The scripts provide reproducibility checks, not proofs of
the infinite statements.
\end{document}
